\section{Introduction}
\label{sec:introduction}

\emph{Equality saturation} mitigates the phase-ordering problem in term
rewriting: applying one rule may enable or disable another, so the order of
application can determine the quality of the final term.  Equality saturation
instead applies rules non-destructively until reaching a fixed point or a
resource limit.  It is powered by \emph{e-graphs}~\cite{nelson1980techniques},
which compactly represent a congruence by grouping equivalent terms into
\emph{e-classes} and sharing their common subterms.

Most work on e-graphs concerns algorithms and applications.  Their semantics
is commonly described through an equality-saturation algorithm~\cite{EggPaper},
a Datalog program~\cite{zhang_relational_2022}, or tree
automata~\cite{tree_automata_e_graphs}.  These accounts explain how saturation
is computed, but a semantics as merely a set of equivalent terms forgets the
sharing and cycles present in the e-graph itself.

Tiurin et al.~\cite{tiurin2025equivalencehypergraphsdporewriting} instead
interpret e-graphs as string diagrams for semilattice-enriched cartesian
categories.  E-classes are represented by joins, sharing by cartesian copying,
and cycles by trace.  
E-graph operations are then interpreted as string-diagram rewriting,
formalised as double-pushout rewriting of hypergraphs.  This account allows
e-graphs to be generalised from cartesian theories to symmetric monoidal
theories with genuinely multi-output generators, such as functions returning
joint probability distributions.  Hypergraphs are particularly useful in
this setting because their isomorphisms absorb many of the structural
equations of symmetric monoidal categories.

We develop the type-theoretic counterpart of this account: a calculus in which
an e-graph is a term and equality saturation is term rewriting.  The
semilattice join $M\join N$ records two members of an e-class; an explicit
\textbf{let} records sharing; and \textbf{letrec} records cyclic sharing.
Unlike a point-free syntax with explicit identities and composition, the
calculus absorbs the unit and associativity laws of composition into variables
and substitution. 

\begin{figure}[t]
  \centering
  \adjustbox{width=\textwidth}{
    \tikzfig{figures/e-graph-example-2}
  }
  \caption{Equality saturation, conventional e-graphs, and their
    string-diagram representations~\cite[Fig.~1]{tiurin2025equivalencehypergraphsdporewriting}.}
  \label{fig:e-graph-example}
\end{figure}

\Cref{fig:e-graph-example} recalls a standard equality-saturation example.
Each column shows a term (or rewrite rule), its e-graph, and the corresponding
string diagram.  An e-class, drawn as a dashed box, represents a set of
equivalent terms.  Its e-nodes---the round boxes within it---record the
possible root operations and refer to other e-classes as children.
Choosing one e-node from each e-class extracts a term.  In column~(b), the rule
$x*2\rightsquigarrow x\mathbin{<\!\!<}1$ creates a non-singleton class while
retaining the shared occurrences of $a$ and $2$.
Each rewrite rule instantiates its variables with matching e-classes.  In
column~(b), for example, $x$ is instantiated with the class containing $a$;
the other classes are matched by their e-nodes.  The right-hand side is then
added using the same instantiation, and its root class is merged with that of
the left-hand side.  In string diagrams, dashed boxes represent the
semilattice structure on hom-sets, while the cartesian copy $\comonoid$ and
delete $\counit$ generators express sharing and discarded inputs.  Every
component of a dashed box consequently has the same context.

The merge from (d) to (e) makes the class containing $a$ recursive.  The
resulting e-graph represents infinitely many finite unfoldings, including
$a$, $a*1$, and $((a*2)/2)*1$.  In the string diagram this feedback is a
trace.  Thus the two features that are invisible in a plain set-of-terms
semantics---sharing and cyclic sharing---are explicit categorical structure.
Extracting a term from an e-graph means choosing an e-node and recursively
following its children.  In a string diagram one instead extracts a component
by rewriting; following its wires may require, for example, sliding for the
trace and naturality of deletion.

Our terms expose the same structure syntactically.  The two most illustrative
columns of \Cref{fig:e-graph-example} become the terms in
\Cref{lst:terms-e-graph}.
\begin{listing*}[t]
\centering
\begin{minipage}[t]{0.43\textwidth}
  \begin{minted}[escapeinside=||]{haskell}
let x = a in
let z = 2 in
let e = x * z |$\join$| x << 1 in
  e / z
  \end{minted}
  \subcaption*{(b) Acyclic sharing}
\end{minipage}
\hfill
\begin{minipage}[t]{0.53\textwidth}
  \begin{minted}[escapeinside=||]{haskell}
let z = 2 in
let w = 1 in
|\kw{letrec}| x =
  (let e = x * z |$\join$| x << w in e / z)
  |$\join$| (let c = z / z |$\join$| w in x * c)
  |$\join$| a
in x
  \end{minted}
  \subcaption*{(e) Cyclic sharing}
\end{minipage}
\caption{Terms for e-graphs (b) and (e) of
  \Cref{fig:e-graph-example}.}
\label{lst:terms-e-graph}
\end{listing*}
In~(b), the binding of $e$ names a whole e-class, so the division consumes one
shared class rather than two duplicated alternatives.  In~(e), the variable
$x$ occurs in its own defining class, and \textbf{letrec} makes that cycle
finite in the syntax.  These examples motivate the two principal extensions
of the core calculus developed below.  We work with representable
semilattice-enriched multicategories, encompassing both the symmetric
monoidal and cartesian settings used by Tiurin et
al.~\cite{tiurin2025equivalencehypergraphsdporewriting}.

\subsection{Contributions}

We make the following contributions.
\begin{itemize}
  \item We give a type theory for representable $\catname{SLat}$-multicategories
  (\Cref{def:s_lat_type_theory}).  Although its distributive rules make
  derivations non-unique, we prove coherence modulo permutation of join
  introduction, define substitution on derivations, and prove derivation
  independence, the substitution laws, and multilinearity
  (\Cref{theorem:coherence},
  \Cref{prop:subst_derivation_independence}, and
  \Cref{lemma:subst_properties}).
  The resulting calculus presents the free representable
  $\catname{SLat}$-multicategory (\Cref{theorem:free_slat_multicat}).

  \item We extend the calculus with \textbf{let} for explicit sharing and
  \textbf{letrec} for cyclic sharing
  (\Cref{def:explicit_sharing,def:recursive_sharing}).  The latter is
  interpreted by trace; because trace need not preserve joins, we develop a
  binding-sensitive coherence and substitution argument and prove the
  corresponding freeness theorem (\Cref{theorem:free_traced_slat}).
\end{itemize}
